% No preamble. Composite macrostep integral, unique prefix cmi-.
\section{The exact composite integral and its pre-gauge continuation}
\label{cmi-section}

Fix the primitive odd ratio $\ell$ and a macro ratio $R=\ell^s$.
The preceding refinement is $M=\ell^k$. All statements concern
finite compatible periodic regulators. Constants may depend on the
fixed $R,\ell$, but are uniform in $M$ and volume. The constraint is
the actual composite
\[
 \mathcal F_R=\mathcal M_\ell^s.
\]
It is not replaced by a direct-$R$ ordinary average. The direct-$R$
internal contour variable $P_R$ is used independently for gauge
fixing. Nor is the one-shot integral below identified with $s$
separately gauge-fixed and truncated quantum integrations.

\subsection{A gauge insertion with its exact normalization}
\label{cmi-gauge-section}

Use normalized Haar measure on the compact matrix group $G$, and
the convention $V^u(b)=u(b_-)V(b)u(b_+)^{-1}$. The previously
proved local geometry gives
\begin{align}
 \mathcal F_R(V^u)
   &=\mathcal F_R(V)^{u|_{\text{final centers}}},\notag\\
 P_R(V^u;y,x)&=u(y)P_R(V;y,x)u(x)^{-1}.                 \label{cmi-covariance}
\end{align}
Every contour defining $P_R$ stays inside its $R$-block.
Put
\begin{equation}\label{cmi-penalty}
 v(h)=1-\operatorname{Re}\tau h,\qquad
 \mathcal G_R(V)=\sum_y\sum_{\substack{x\in B_R(y)\\x\ne y}}
                       v(P_R(V;y,x)).
\end{equation}
The analytic continuation of $v(h)$ is
$1-\tfrac12\tau(h+h^{-1})$. This is the absolute contour
penalty. It is not replaced by a relative penalty away from
a gauge-zero background.

A local logarithmic formula does not by itself allow integration
over every compact gauge transformation. One may work on a
gauge-invariant domain where the contour formula is defined on
complete orbits, or make the following explicit measurable
completion. For a path tuple $(U_j)$ from $y$ to $x$, retain its
specified reference $T$. On a fixed small conjugation-invariant
neighborhood of the identity for every $U_jT^{-1}$, put
\[
 P(U_1,\ldots,U_n)
   =\Phi(U_1T^{-1},\ldots,U_nT^{-1})T,
\]
where $\Phi$ is the \emph{actual selected local source mean}
on the reference-normalized tuple, with its original contour
multiplicities, ordering and reference. Right endpoint covariance
proves this representation on the original domain, and simultaneous
conjugation covariance follows from the same source covariance.
On the complementary set put $P=T$. Thus the completion is
measurable, group valued and endpoint covariant everywhere,
and equals the chosen analytic mean near the identity. No
different mean or additional mean equation is substituted.

Similarly, for each primitive ordinary average, use its original
logarithmic formula when all its reference-normalized transported
strings lie in the chosen small logarithm neighborhood, and its
reference string otherwise. The condition is conjugation invariant.
This gives a measurable gauge-covariant completion of
$\mathcal M_\ell$, hence of its composition $\mathcal F_R$.
The values outside the analytic domain are part of this proposed
full blocking map; equality with a different large-field completion
is not claimed. The following local formulas agree wherever
the original analytic maps agree.

Let $\mathcal H_R$ be the group of current-site gauges equal to
one at final block centers, with normalized independent Haar
measure at the $n_s$ other sites. For $g>0$ define
\begin{equation}\label{cmi-gauge-constant}
 c(g)=\int_G e^{-v(h)/g^2}\,dh>0.
\end{equation}
Then
\begin{equation}\label{cmi-orbit-identity}
 \int_{\mathcal H_R}e^{-\mathcal G_R(V^u)/g^2}\,du
       =c(g)^{n_s}.
\end{equation}
Indeed, by \eqref{cmi-covariance}, the factor at $(y,x)$ is
$e^{-v(P_R(V;y,x)u(x)^{-1})/g^2}$, whose integral is $c(g)$
by Haar invariance, independently of $V$. The site integrals
factor. There is no omitted gauge determinant: this is direct
integration over compact Haar variables with the displayed
normalization.

For an integrable gauge-invariant current density $\rho$, and
any bounded test function $f$ of the final configuration,
\begin{equation}\label{cmi-pushforward}
 \int dV\,\rho(V)f(\mathcal F_R(V))
 =c(g)^{-n_s}\int dV\,\rho(V)f(\mathcal F_R(V))
                            e^{-\mathcal G_R(V)/g^2}.
\end{equation}
Insert \eqref{cmi-orbit-identity}, use Fubini, then change
$V$ by the fixed gauge $u$; link Haar, $\rho$ and $\mathcal F_R$
are unchanged. This is an identity of finite pushforward
measures, requiring no globally regular delta density.
It also holds for a finite gauge-invariant current measure,
and for a gauge-invariant orbit-domain restriction.

A noninvariant small-field window is introduced \emph{after}
this insertion. Splitting the gauged integral into $\chi_{\rm sf}$
and $1-\chi_{\rm sf}$ is exact for any measurable
$0\le\chi_{\rm sf}\le1$ on the fiber. If $\rho$ is positive
both pieces are positive on real fields. Neither piece is
asserted to satisfy \eqref{cmi-pushforward} without the gauge
factor. The complement requires a separate large-field estimate.

\subsection{Identification of the physical centered branch}
\label{cmi-branch-section}

Let $U_f(W)$ be the proved full ordinary minimizer at total
refinement $MR$, with final spacing $R$ in current units, and set
\begin{equation}\label{cmi-pregauge}
 V_p(W)=\mathcal M_\ell^k(U_f(W)),\qquad
 \mathcal F_R(V_p(W))=W.
\end{equation}
Let $U_c(V)$ be the preceding-depth ordinary minimizer and
$S_M(V)=S(U_c(V))$. Choose a common current comparison radius
$\delta_R$ and then a final datum radius $b_R$ so that
$U_f(W)$ belongs to the current minimum class, every $U_c(V)$
with $\|\log V/i\|_\infty<\delta_R$ belongs to the full
final-spacing-$R$ minimum class, and the partial average and
gauges below stay within all common chart and invariance domains.
These are finitely many strict smallness conditions with
$R$-dependent constants, uniform in $M$ and volume. In dimension
four the coupling-free Wilson action has no additional scale
factor; the potential norms acquire only fixed powers of $R$.

The two actual comparisons are
\begin{equation}\label{cmi-two-comparisons}
 S_M(V_p)\le S(U_f)
       \le S(U_c(V_p))=S_M(V_p).
\end{equation}
For the second inequality exact composition makes $U_c(V_p)$
a full feasible competitor. For any other current $V$ in the
comparison ball with $\mathcal F_R(V)=W$, the same argument gives
\begin{equation}\label{cmi-current-minimum}
 S_M(V)\ge S(U_f)=S_M(V_p).
\end{equation}
This is a common-domain minimum comparison, not a
dynamic-programming assertion for arbitrary backgrounds.

Define the root-fixed current gauge and background by
\begin{equation}\label{cmi-rooted-background}
 u(y)=1,\qquad u(x)=P_R(V_p;y,x),\qquad V_*(W)=V_p(W)^u.
\end{equation}
Then $\mathcal F_R(V_*)=W$ and $P_R(V_*;y,x)=1$ exactly.
At fixed $R$ this is local analytic with
$\|\log u/i\|_\infty\le C_R\|\log V_p/i\|_\infty$.
The already proved coarse-action gauge invariance, with its
smooth fine point-value extension, implies $S_M(V_*)=S_M(V_p)$
after shrinking the common radii. Since $\mathcal G_R\ge0$
and $\mathcal G_R(V_*)=0$ on real fields, $V_*$ is a local
constrained minimum of $S_M+\mathcal G_R$.

The preceding Gaussian module supplies its uniformly positive
flat retained precision. The local analytic chart and
$D^2S_M=q_M+O(\|\log V\|_\infty)$ show that the pulled-back
Hessian on a common convex retained ball is
$A_{M,0}+O_R(\|w\|_\infty+\|z\|_\infty)$. It remains positive
after a fixed shrink. Therefore the centered real critical
point is unique in that common neighborhood; gauge zero
alone would not prove uniqueness. Constraint, gauge-zero and
first-derivative identities continue holomorphically to complex
$W$. No complex minimum is asserted.

The supported pre-gauge theorem gives the stronger amplitude
bound
\begin{equation}\label{cmi-pregauge-amplitude}
 a_p=\log V_p/i,\qquad
 \|a_p\|_\infty\le C_a\|w\|_\infty/R.
\end{equation}
Its supported-family constant is fixed at primitive $\ell$
independently of $s,k$. There is no corresponding $1/R$
contraction claim for the rooted potential $\log V_*/i$.

\subsection{The local density, Haar factors and pivot determinants}
\label{cmi-density-section}

Use Hermitian Lie coordinates orthonormal for $\tau$, of
dimension $q=\dim\mathfrak g$. Let $n_e$ and $n_c$ count
positive current and final macro bonds, and put
$n_z=q(n_e-n_c)$. In exponential coordinates write normalized
Haar as
\begin{equation}\label{cmi-Haar}
 d(e^{iX})=h_0j(X)\,dX,\qquad h_0>0,\quad j(0)=1.
\end{equation}
Consequently
$\delta_G(e^{iF})=h_0^{-1}\delta_{\mathbb R^q}(F)$,
with the coordinate measure understood: $j(F)$ is evaluated
at $F=0$.

Fix a real small $W$ and $V=V_*(W)$. The composite relative
constraint and its pivot derivative are
\begin{equation}\label{cmi-relative}
 F_V(B)=\log[\mathcal F_R(e^{iB}V)\mathcal F_R(V)^{-1}]/i,
 \qquad A_V=D F_V(0),\qquad
 A_{c,p(c)}=R^{1-d}\pi_c(V).
\end{equation}
Let $\psi_V$ be the ambient straightening, $K_V$ the linear
retained insertion, and $\Phi_V(z)=\psi_V(K_Vz)$.
Each retained component of $\Phi_V$ equals $z$ exactly.
Order ambient variables as retained coordinates followed by
pivots. The derivative of $\psi_V$ is block triangular;
its pivot diagonal blocks satisfy
\[
 D_{p(c)}F_V(\psi_V(B);c)\,
 D_{B(p(c))}\psi_V(B)(p(c))=A_{c,p(c)}.
\]
Thus
\begin{equation}\label{cmi-ambient-Jacobian}
 \det D\psi_V(B)=
   \prod_c\frac{\det A_{c,p(c)}}
                  {\det D_{p(c)}F_V(\psi_V(B);c)}.
\end{equation}
Let $j_c^{\rm str}(V;B)$ be the analytic logarithm of the
$c$th ratio, zero at $B=0$, and put
\begin{equation}\label{cmi-background-delta}
 \mathcal D_R(V)=\prod_c(\det A_{c,p(c)})^{-1}.
\end{equation}
The determinants have their positive flat orientation on the
real small chart.

Choose the small-field sector to be the component of this
relative chart with retained coordinates in a specified product
window, retaining any additional indicator explicitly. Changing
the square ambient variables and then solving $A_VB=0$ gives
its exact density
\begin{align}
 &c(g)^{-n_s}h_0^{n_e-n_c}\mathcal D_R(V)
 \int dz\,
    \prod_bj(\Phi_V(z)(b))\,
    e^{\sum_cj_c^{\rm str}(V;K_Vz)}
    \rho(e^{i\Phi_V(z)}V)
    e^{-\mathcal G_R(e^{i\Phi_V(z)}V)/g^2}.
                                                        \label{cmi-local-density}
\end{align}
The integration domain and every extra indicator are those
of the chosen sector. This uses the square Jacobian
\eqref{cmi-ambient-Jacobian}, not a determinant of the
rectangular derivative $D\Phi_V$. Equivalently the unrearranged
pivot density is the product of inverse
$D_{p(c)}F_V(\Phi_V(z);c)$ determinants.

Here the globally gauge-invariant $\rho$ in
\eqref{cmi-pushforward} is an input density or measure.
Only on the common small chart do we use a representation
\[
 \rho(U)=\exp[-S_M(U)/g^2+E(U)],
\]
with any field-independent factor tracked separately.
The local classical branch $S_M$ has not been promoted
to a global action on all configurations.

Put $z=g\xi$ for real $g>0$, and let $A=A_M(V)$ be the exact
retained Hessian of $S_M+\mathcal G_R$. Define
\begin{align}
 \mathcal R_3(V,z)
   &=S_M(e^{i\Phi_V(z)}V)+\mathcal G_R(e^{i\Phi_V(z)}V)
            -S_M(V)-\tfrac12\langle z,Az\rangle,\notag\\
 I(V,g,\xi)
   &=-g^{-2}\mathcal R_3(V,g\xi)
      +\sum_cj_c^{\rm str}(V;K_Vg\xi)
      +\sum_b\log j(\Phi_V(g\xi)(b))\notag\\
   &\hspace{1cm}+E(e^{i\Phi_V(g\xi)}V)-E(V).
                                                        \label{cmi-interaction}
\end{align}
There is no linear term because the background is the physical
centered branch. Formula \eqref{cmi-local-density} becomes
\begin{align}
 &c(g)^{-n_s}h_0^{n_e-n_c}\mathcal D_R(V)\,
   g^{n_z}(2\pi)^{n_z/2}\det(A)^{-1/2}
   e^{-S_M(V)/g^2+E(V)}\notag\\
 &\hspace{3cm}\times
   \mathbb E_{A^{-1}}\!
      \left[\chi_{\rm window}(\xi)e^{I(V,g,\xi)}\right].
                                                        \label{cmi-Gaussian-integral}
\end{align}
The Gaussian convention is the retained Euclidean density
$\det(A)^{1/2}(2\pi)^{-n_z/2}
e^{-\langle\xi,A\xi\rangle/2}$.
The covariance is that of this exact classical
action-plus-gauge Hessian; Haar remains in
\eqref{cmi-interaction}. The factors $c(g)^{-n_s}$,
$h_0^{n_e-n_c}$ and $g^{n_z}$ are vacuum normalizations.
$\mathcal D_R(V)$ and $\det A$ generally depend on the
background and are not vacuum constants.

\subsection{The literal retained cutoff}
\label{cmi-window-section}

For the original retained link/color balls set
\[
 \chi_{\rm window}(\xi)=
 \prod_{b\ {\rm retained}}\mathbf1_{\{|\xi(b)|\le T\}}.
\]
Since $\Phi_V(g\xi)(b)=g\xi(b)$ at every retained bond,
this is exactly the relative-coordinate window of the
chosen sector. It is not a cutoff on $K_V\xi$, or a
Euclidean ball on a whole macro fiber. Pivot, contour,
plaquette or reconstructed-field indicators, if included,
remain at their exact $\Phi_V(g\xi)$ pullbacks.
The macro sector is not identified with an intersection
of the primitive windows of $s$ separate integrations.

At fixed $R$, the normalized pivot logarithms are local
$O_R(|g|T)$ terms. Inversion invariance gives
$\log j(X)=O(|X|^2)$, so Haar logarithms are local
$O_R(g^2T^2)$ terms. The gauge third-and-higher remainder
is local $O_R(|g|T^3)$. The internal contours read no pivot,
and therefore its full exact factor satisfies
\begin{equation}\label{cmi-gauge-retained}
 \mathcal G_R(e^{i\Phi_V(z)}V)
       =\mathcal G_R(e^{iE_{\rm free}z}V).
\end{equation}
On an independent background product ball its constant
and linear terms must also be separated away from the
physical gauge-zero branch. All these estimates have
finite local supports and finite anchor multiplicities
at fixed $R$, without a sum over the old refinement $M$.
The classical action and inherited interaction require
their own support estimates.

\subsection{Pre-gauge continuation of the full inherited insertion}
\label{cmi-continuation-section}

With $V_*=V_p^u$, set, link by link,
\begin{equation}\label{cmi-prefluctuation}
 B_p(w,\xi)(b)=
   \operatorname{Ad}_{u(w;b_-)^{-1}}
                  \Phi_{V_*(w)}(g\xi)(b).
\end{equation}
Direct multiplication proves
\begin{equation}\label{cmi-fluctuation-identity}
 \bigl(e^{i\Phi_{V_*}(g\xi)}V_*\bigr)^{u^{-1}}
             =e^{iB_p(w,\xi)}V_p(w).
\end{equation}
The coupling remains inside the nonlinear chart:
$\Phi_V(g\xi)$ is not replaced by $g\Phi_V(\xi)$.
The gauge depends on $w$, not on $\xi$. Its real adjoint
is a color isometry; on the small complex gauge ball it
has a fixed $R$-dependent bound.

Let $E$ be a holomorphic total inherited functional on a
current potential product ball of radius $b$, with its
local gauge invariance on the original common domain.
Individual activity coefficients need not be gauge invariant.
On a sufficiently small common overlap, both arguments of
\eqref{cmi-fluctuation-identity}, and a small connecting
gauge path, are in this domain. Their $E$ values therefore
agree, first on real fields and then by holomorphic uniqueness.
If only $e^E$ is supplied, a compatible nonvanishing
holomorphic logarithm branch must be supplied before making
the corresponding statement about $E$.

On a fixed target $b$-ball below all classical domains,
\eqref{cmi-pregauge-amplitude} and the local chart give
\begin{equation}\label{cmi-common-ball}
 \left\|\frac1i\log[e^{iB_p(w,\xi)}V_p(w)]\right\|_\infty
       \le C_a b/R+C_R|g|T.
\end{equation}
Choose $R\ge4C_a$ and $C_R|g|T\le b/4$.
The pre-gauge argument then stays inside the old $b$-ball
uniformly in $M$ and volume. Its inherited functional is
defined and holomorphic on the whole target ball, even
where the rooted argument falls outside the original
domain of $E$. It is a continuation of the common-domain
germ, not an assertion that $E$ was already defined at
that rooted argument. Replacing $g\xi$ by $\lambda g\xi$
gives a complex $\lambda$-disc of radius at least
$cb/(C_R|g|T)$ within the strict chart margins.

The \emph{entire} inherited contribution in
\eqref{cmi-Gaussian-integral} is continued: replace its
outside $E(V_*)$ by $E(V_p)$, and replace the inside
increment by
\[
 E(e^{iB_p(w,\xi)}V_p(w))-E(V_p(w)).
\]
Leaving $E(V_*)$ outside would be undefined on part of
the enlarged domain. There is no change of integration
variable $\xi$. The Gaussian, its determinant, the pivot
and Haar factors, the absolute gauge penalty and the
retained cutoff remain exactly those of
\eqref{cmi-Gaussian-integral}; no adjoint Jacobian or
transformed window is introduced.

\subsection{What holomorphy does and does not imply}
\label{cmi-scope-section}

At fixed finite regulator and fixed real retained window,
the continued integrand is holomorphic in $w$. The proved
complex precision neighborhood supplies a consistent
analytic determinant square root. Every compact subdomain
has an integrable uniform majorant, so the fixed-window
integral is holomorphic by differentiation under the
integral. Equality on the nonempty overlap identifies it
as the continuation of the original small-field germ.

For real data and positive physical input, the exact
integrand is positive and the retained window has positive
measure. The integral is then strictly positive. This
does not imply nonvanishing on a volume-independent
complex ball: an integral of pointwise nonzero analytic
factors can vanish. A finite-volume logarithm germ is
available near a positive point; a uniform complex
logarithm needs the full cluster hypotheses and budgets.

Similarly, pointwise $\lambda$-Cauchy applied to a total
$E$ can give only an extensive estimate. An anchored
interaction bound needs a strict joint $(w,\xi)$ family
for the pre-gauge argument, with remote background
supports and an explicit local fluctuation collar.
That separate support module is not inferred from
\eqref{cmi-common-ball}. Purely classical coefficients
have independent fixed classical domains; the quantum
radius may be chosen below them once, rather than
reused as their outer analytic radius at every step.

The present module proves the finite-regulator
normalization and local density identities, identifies
the physical centered branch, preserves the exact
window, and supplies the pointwise pre-gauge extension.
It does not establish a stable quantum iteration,
uniform nonvanishing without the cluster bounds,
large-field control, mixed-history matching, coupling
flow or a continuum measure.
